% ---
% Reemplaza líneas 1037-1092 de Docs/Technical_Report.tex
% Decisiones: N/A (sin cambio)
% ---

\subsubsection{Compuertas cuánticas de un qubit}
% -----------------------------------------------------------

Las operaciones sobre qubits se realizan mediante \textit{compuertas cuánticas}, representadas por matrices unitarias $U \in \mathbb{C}^{2\times 2}$ que satisfacen $UU^\dagger = I$. La unitariedad garantiza la reversibilidad de la operación y la conservación de la norma del estado. Las compuertas más relevantes para los circuitos variacionales son las rotaciones parametrizadas alrededor de los ejes de la esfera de Bloch \cite{nielsen2010quantum}:

\begin{equation}
    R_X(\theta) = e^{-i\frac{\theta}{2}X} = 
    \begin{pmatrix} 
        \cos\frac{\theta}{2}  & -i\sin\frac{\theta}{2} \\ 
        -i\sin\frac{\theta}{2} & \cos\frac{\theta}{2} 
    \end{pmatrix}
    \label{eq:rx}
\end{equation}

\begin{equation}
    R_Y(\theta) = e^{-i\frac{\theta}{2}Y} = 
    \begin{pmatrix} 
        \cos\frac{\theta}{2}  & -\sin\frac{\theta}{2} \\ 
        \sin\frac{\theta}{2}  & \cos\frac{\theta}{2} 
    \end{pmatrix}
    \label{eq:ry}
\end{equation}

\begin{equation}
    R_Z(\theta) = e^{-i\frac{\theta}{2}Z} = 
    \begin{pmatrix} 
        e^{-i\frac{\theta}{2}} & 0 \\ 
        0 & e^{i\frac{\theta}{2}} 
    \end{pmatrix}
    \label{eq:rz}
\end{equation}

\noindent donde $X$, $Y$ y $Z$ son las matrices de Pauli:

\begin{equation}
    X = \begin{pmatrix} 0 & 1 \\ 1 & 0 \end{pmatrix}, \quad
    Y = \begin{pmatrix} 0 & -i \\ i & 0 \end{pmatrix}, \quad
    Z = \begin{pmatrix} 1 & 0 \\ 0 & -1 \end{pmatrix}
    \label{eq:pauli}
\end{equation}

Estas compuertas de rotación son los bloques de construcción del ansatz variacional implementado en Qiskit, donde los ángulos $\theta$ son los parámetros entrenables del circuito. La compuerta de Hadamard, definida como:

\begin{equation}
    H = \frac{1}{\sqrt{2}}\begin{pmatrix} 1 & 1 \\ 1 & -1 
    \end{pmatrix}
    \label{eq:hadamard}
\end{equation}

\noindent transforma el estado base $|0\rangle$ en la 
superposición uniforme 
$\frac{1}{\sqrt{2}}(|0\rangle + |1\rangle)$, y es 
ampliamente utilizada como compuerta de inicialización en 
circuitos cuánticos.
